\documentclass[9pt,a4paper]{extarticle}

\usepackage[a4paper, top=14mm, bottom=14mm, left=14mm, right=14mm]{geometry}
\usepackage[T1]{fontenc}
\usepackage{lmodern}
\usepackage{microtype}
\usepackage[table]{xcolor}
\usepackage{booktabs}
\usepackage{tabularx}
\usepackage{array}
\usepackage{parskip}
\usepackage{enumitem}
\usepackage{titlesec}
\usepackage{mdframed}
\usepackage{multicol}
\usepackage{fancyhdr}
\usepackage{hyperref}
\usepackage{fontawesome5}

% ── Colours ───────────────────────────────────────────────────────────────────
\definecolor{DG}{HTML}{1B5E20}
\definecolor{MG}{HTML}{2E7D32}
\definecolor{LG}{HTML}{43A047}
\definecolor{PG}{HTML}{F1F8E9}
\definecolor{CG}{HTML}{C8E6C9}
\definecolor{TEA}{HTML}{00696E}
\definecolor{AMB}{HTML}{E65C00}
\definecolor{RED}{HTML}{C62828}
\definecolor{RLT}{HTML}{FFEBEE}
\definecolor{ALT}{HTML}{FFF3E0}
\definecolor{SLT}{HTML}{37474F}
\definecolor{INK}{HTML}{212121}
\definecolor{MUT}{HTML}{78909C}

% ── Layout ────────────────────────────────────────────────────────────────────
\setlength{\parskip}{2pt}
\setlength{\columnsep}{8pt}
\setlist[itemize]{leftmargin=1.2em, itemsep=0pt, parsep=0pt, topsep=1pt, label=\textcolor{MG}{\textbullet}}
\setlist[enumerate]{leftmargin=1.4em, itemsep=0pt, parsep=0pt, topsep=1pt}

\titleformat{\section}
  {\bfseries\small\color{white}}
  {}
  {0em}
  {\setlength\fboxsep{3pt}\colorbox{DG}}
\titlespacing{\section}{0pt}{6pt}{3pt}

% ── Header ────────────────────────────────────────────────────────────────────
\pagestyle{fancy}
\fancyhf{}
\renewcommand{\headrulewidth}{0pt}
\fancyhead[L]{\tiny\color{MUT} 300,000 Streets of Melbourne --- Quick Revision Sheet}
\fancyhead[R]{\tiny\color{MUT} Regen Melbourne \texttimes{} RMIT \textbullet{} Hishikesh Phukan}
\fancyfoot[C]{\tiny\color{MUT}\thepage}

% ── Helpers ───────────────────────────────────────────────────────────────────
\newcommand{\kw}[1]{\textbf{\color{DG}#1}}
\newcommand{\num}[1]{\textbf{\color{RED}#1}}
\newcommand{\good}[1]{\textbf{\color{MG}#1}}
\newcommand{\warn}[1]{\textbf{\color{AMB}#1}}
\newcommand{\bad}[1]{\textbf{\color{RED}#1}}
\newcommand{\tag}[2]{\colorbox{#1}{\tiny\color{white}\textbf{#2}}}

\newcolumntype{L}[1]{>{\raggedright\arraybackslash}p{#1}}
\newcolumntype{C}[1]{>{\centering\arraybackslash}p{#1}}

% =============================================================================
\begin{document}

% ── Title strip ───────────────────────────────────────────────────────────────
\colorbox{DG}{\parbox{\linewidth}{%
  \vspace{3pt}\hspace{6pt}%
  {\color{white}\large\bfseries 300,000 Streets of Melbourne}%
  \hfill{\color{LG}\small School Streets Safety --- Quick Revision Sheet}%
  \hspace{6pt}\vspace{3pt}%
}}
\vspace{4pt}

% =============================================================================
\begin{multicols}{2}

% ── THE ONE-LINE PITCH ────────────────────────────────────────────────────────
\section{The pitch}
A reproducible Python pipeline that \kw{scores}, \kw{maps}, and \kw{predicts} the safety of school walking routes using open government data and the \kw{Healthy Streets} framework --- then surfaces costed fixes for councils.

% ── THE NUMBERS THAT MATTER ───────────────────────────────────────────────────
\section{Numbers to remember}

\begin{tabular}{@{}L{1.6cm} L{5.6cm}@{}}
\num{192\%}   & crash increase in Darebin LGA, 2021\,\textrightarrow\,2024 (25\,\textrightarrow\,73) \\
\num{0.4}     & Preston HS HS2 score --- no safe crossing \\
\num{r=0.84}  & Pearson correlation: SEIFA decile vs HS score \\
\num{100\%}   & William Ruthven SC crashes during school hours \\
\num{7,948}   & Victorian crash records analysed (2021--2025) \\
\num{26}      & open-data features fed into ML model \\
\num{MAE 2.88}& Ridge regression mean absolute error (LOO-CV, n=3) \\
\num{30}      & pre-computed scenario results (10 templates \texttimes{} 3 schools) \\
\num{\$80k}   & cheapest fix for Preston HS (signalised crossing) \\
\num{200+}    & Melbourne secondary schools scalable to with no code changes \\
\end{tabular}

% ── 10 HS INDICATORS ─────────────────────────────────────────────────────────
\section{10 Healthy Streets indicators (threshold: 6.0 / 10)}

\begin{tabular}{@{}C{0.56cm} L{1.8cm} L{1.6cm} L{2.6cm}@{}}
\rowcolor{DG}
\color{white}\tiny\textbf{Code} &
\color{white}\tiny\textbf{Indicator} &
\color{white}\tiny\textbf{Source} &
\color{white}\tiny\textbf{Worst score} \\
\rowcolor{PG}
HS1 & Footpath      & Field   & Reservoir 4.2 \warn{[MOD]} \\
HS2 & Easy to cross & Field   & \bad{Preston 0.4 [MAJOR]} \\
\rowcolor{PG}
HS3 & Shade         & OSM     & W.Ruthven 3.0 \warn{[MOD]} \\
HS4 & Rest places   & OSM     & W.Ruthven 3.3 \warn{[MOD]} \\
\rowcolor{PG}
HS5 & Not too noisy & Field   & Reservoir 5.0 \\
HS6 & Active travel & OSM+fld & W.Ruthven 4.2 \\
\rowcolor{PG}
HS7 & Feel safe     & CSA     & \good{All $\geq$ 7.5} \\
HS8 & Things to do  & OSM     & W.Ruthven 4.3 \\
\rowcolor{PG}
HS9 & Feel relaxed  & Field   & Reservoir 5.2 \\
HS10& Clean air     & EPA     & \good{All $\geq$ 9.0} \\
\end{tabular}

\vspace{3pt}
\kw{Severity rules:} \bad{Major} if HS2\,$<$\,4 OR HS1\,$<$\,4 OR HS5\,$<$\,3 \textbullet{}
\warn{Moderate} if 2+ indicators\,$<$\,6.0

% ── 3 SCHOOLS AT A GLANCE ────────────────────────────────────────────────────
\section{3 schools at a glance}

\begin{tabular}{@{}L{2.2cm} C{0.9cm} C{1.0cm} L{3.2cm}@{}}
\rowcolor{DG}
\color{white}\tiny\textbf{School} &
\color{white}\tiny\textbf{Overall} &
\color{white}\tiny\textbf{Severity} &
\color{white}\tiny\textbf{Key issue} \\
Preston HS     & \bad{7.2} & \tag{RED}{MAJOR}    & HS2=0.4 no crossing \\
\rowcolor{PG}
Reservoir HS   & \warn{6.1}& \tag{AMB}{MOD}      & HS1=4.2 broken footpath \\
W. Ruthven SC  & \warn{6.7}& \tag{AMB}{MOD}      & HS3=3.0 no shade \\
\end{tabular}

\vspace{2pt}
\textit{\small The paradox: Preston has the \bad{highest} overall score but the \bad{only Major} severity. One failing indicator overrides nine green ones.}

% ── PIPELINE IN 10 STEPS ─────────────────────────────────────────────────────
\section{Pipeline --- 10 steps, 1 command}

\texttt{\small python run\_all.py}

\begin{enumerate}
  \item \kw{Ingest} --- crash (VicRoads), streets (OSM), air (EPA), crime (CSA), SEIFA, Census
  \item \kw{Feature engineering} --- 26 features per school gate; filter sentinels 777/888/999
  \item \kw{HS scoring} --- \texttt{src/scoring/hs1.py} \ldots \texttt{hs10.py}; each scorer is independent
  \item \kw{Severity classify} --- rule-based thresholds in \texttt{severity.py}
  \item \kw{Recommendations} --- rule engine: indicator gap \textrightarrow{} intervention + cost + timeline
  \item \kw{ML training} --- Ridge regression, LOO-CV, 26 features, predict HS1--HS10
  \item \kw{Scenario engine} --- 10 templates \texttimes{} 3 schools = 30 delta results
  \item \kw{Equity \& crash analysis} --- SEIFA join, 400\,m crash buffers, hourly breakdown
  \item \kw{Visualisation} --- 9 charts (Matplotlib); GeoPackage + QGIS project export
  \item \kw{Dashboard export} --- JSON \textrightarrow{} static GitHub Pages (Leaflet + Chart.js)
\end{enumerate}

\columnbreak

% ── ML IN PLAIN TERMS ────────────────────────────────────────────────────────
\section{Machine learning --- what and why}

\begin{itemize}
  \item \kw{Algorithm:} Ridge regression (handles collinear OSM features)
  \item \kw{Features:} 26 (20 spatial via OSM, 2 environmental EPA/CSA, 4 crash stats)
  \item \kw{Evaluation:} Leave-One-Out CV --- only valid method with n=3 schools
  \item \kw{Mean MAE:} 2.88 across all 10 indicators
  \item \good{Can automate:} HS7 (crime data) and HS10 (EPA air quality) --- no surveyor needed
  \item \bad{Cannot automate:} HS2 --- crossing \emph{presence} $\neq$ crossing \emph{quality}
\end{itemize}
\vspace{2pt}
\textit{\small Implication: councils can triage all 200+ Melbourne schools on HS7 and HS10 at zero survey cost today.}

% ── EQUITY FINDING ───────────────────────────────────────────────────────────
\section{Equity finding}

\begin{tabular}{@{}L{2.2cm} C{0.8cm} C{0.7cm} L{3.2cm}@{}}
\rowcolor{DG}
\color{white}\tiny\textbf{School} &
\color{white}\tiny\textbf{SEIFA} &
\color{white}\tiny\textbf{HS} &
\color{white}\tiny\textbf{Walking dependency} \\
Reservoir HS   & \bad{3.7}  & \warn{6.1} & High --- lower income, less cars \\
\rowcolor{PG}
W.Ruthven SC   & \warn{3.7} & \warn{6.7} & High --- similar catchment \\
Preston HS     & \good{5.5} & \good{7.2} & Lower --- more car ownership \\
\end{tabular}

\vspace{2pt}
\kw{r\,=\,0.84:} lower decile (more disadvantage) \textrightarrow{} worse safety score.
Lower income = less car access = more walking = more exposure to dangerous streets.

% ── SCENARIO ENGINE ──────────────────────────────────────────────────────────
\section{Scenario engine --- top 3 fixes}

\begin{tabular}{@{}L{2.1cm} L{2.2cm} C{0.6cm} L{2.4cm}@{}}
\rowcolor{DG}
\color{white}\tiny\textbf{School} &
\color{white}\tiny\textbf{Intervention} &
\color{white}\tiny\textbf{Cost} &
\color{white}\tiny\textbf{Impact} \\
Preston HS     & Signalised crossing    & \$80--200k & HS2: 0.4\,\textrightarrow\,3.5; MAJOR\,\textrightarrow\,MOD \\
\rowcolor{PG}
Reservoir HS   & Footpath repair        & \$100--300k& HS1: 4.2\,\textrightarrow\,7.0 \\
W.Ruthven SC   & 40\,km/h school zone   & \$5--20k   & HS9: instant gain \\
\end{tabular}

% ── CODEBASE MAP ─────────────────────────────────────────────────────────────
\section{Codebase map}

\begin{tabular}{@{}L{2.8cm} L{4.5cm}@{}}
\texttt{src/scoring/}        & hs1.py \ldots hs10.py + severity.py \\
\rowcolor{PG}
\texttt{src/features/}       & engineer.py (26 features) \\
\texttt{src/ml/}             & train / predict / evaluate \\
\rowcolor{PG}
\texttt{src/scenarios/}      & 10 intervention delta templates \\
\texttt{src/recommendations/}& rule-to-fix mapping + cost metadata \\
\rowcolor{PG}
\texttt{src/data/}           & download + cache handlers \\
\texttt{src/visualisation/}  & 9 chart builders \\
\rowcolor{PG}
\texttt{docs/}               & GitHub Pages dashboard \\
\texttt{outputs/}            & CSV, PNG, GPKG, HTML, PPTX \\
\end{tabular}

% ── TECH STACK ───────────────────────────────────────────────────────────────
\section{Tech stack (one line)}

Python 3.12 \textbullet{} pandas / geopandas / OSMnx \textbullet{} scikit-learn \textbullet{} Matplotlib \textbullet{} Leaflet.js \textbullet{} Chart.js \textbullet{} GitHub Pages \textbullet{} \LaTeX{}

% ── LIKELY QUESTIONS ─────────────────────────────────────────────────────────
\section{Likely Q\&A}

\begin{itemize}
  \item \kw{Why Ridge?} Collinear OSM features; Ridge handles this; interpretable weights.
  \item \kw{Why LOO-CV?} Only 3 schools --- any train/test split would give 0 test samples.
  \item \kw{Why not random forest?} N=3 gives insufficient variance for ensemble methods.
  \item \kw{How does it scale?} Same \texttt{run\_all.py}; swap school list in \texttt{config.py}; no code changes.
  \item \kw{Is the r=0.84 significant?} With n=3 it's directional not inferential --- acknowledged in analysis.
  \item \kw{Why not HS2 automated?} OSM tags crossing presence; it cannot tag signal timing, visibility, sightlines, or informal crossings.
  \item \kw{What does SEIFA measure?} Relative socio-economic disadvantage (ABS 2021, SA1 level, pop-weighted IRSD).
  \item \kw{Delta method for scenarios?} Anchors to real observed scores; model estimates the \emph{change}, not the absolute value --- avoids compounding prediction error.
\end{itemize}

\end{multicols}

% ── Bottom strip ──────────────────────────────────────────────────────────────
\vspace{4pt}
\colorbox{DG}{\parbox{\linewidth}{%
  \vspace{3pt}\hspace{6pt}%
  {\color{LG}\small\textbf{Dashboard:} s4111078-sidd.github.io/300-000-School-Streets}%
  \hfill%
  {\color{white}\small\textbf{Run everything:} \texttt{python run\_all.py}}%
  \hspace{6pt}\vspace{3pt}%
}}

\end{document}
